


All login endpoints are **unauthenticated**. On success, you receive a `token` that must be sent as `Authorization: Bearer <token>` on every subsequent request.

<Warning>
  Login endpoints are rate-limited per IP. Repeated failed attempts will result
  in a `429 Too Many Requests` response.
</Warning>

---

\subsection{General login}

`POST /auth/login`

Authenticates any user — student or admin — and returns a JWT with the caller's role. Use the role-specific endpoints (`/auth/login/student`, `/auth/login/admin`) when you need to restrict which account types are accepted.

**Auth required:** No

\subsubsection{Request body}

<ParamField body="username" type="string" required>
  The student roll number (e.g. `O210008`) or admin username (e.g. `director`).
</ParamField>

<ParamField body="password" type="string" required>
  The account password.
</ParamField>

\subsubsection{Response}

<ResponseField name="success" type="boolean" required>
  `true` on successful authentication.
</ResponseField>

<ResponseField name="token" type="string" required>
  Signed JWT. Pass this value in the `Authorization: Bearer` header for all
  authenticated requests.
</ResponseField>

<ResponseField name="role" type="string" required>
  The role assigned to this account. Possible values: `student`, `director`,
  `dean`, `warden_male`, `warden_female`, `caretaker_male`, `caretaker_female`,
  `security`, `webmaster`.
</ResponseField>

<ResponseField name="username" type="string" required>
  The username that was authenticated.
</ResponseField>

\subsubsection{Example}

\begin{lstlisting}[language=bash]
curl --request POST \
  --url https://api.uniz.rguktong.in/api/v1/auth/login \
  --header 'Content-Type: application/json' \
  --data '{
    "username": "O210008",
    "password": "password123"
  }'

\end{lstlisting}

**200 OK — student**

\begin{lstlisting}[language=json]
{
  "success": true,
  "token": "eyJhbGciOiJIUzI1NiIsInR5cCI6IkpXVCJ9...",
  "role": "student",
  "username": "O210008"
}

\end{lstlisting}

**200 OK — admin**

\begin{lstlisting}[language=json]
{
  "success": true,
  "token": "eyJhbGciOiJIUzI1NiIsInR5cCI6IkpXVCJ9...",
  "role": "director",
  "username": "director"
}

\end{lstlisting}

\subsubsection{Error responses}

| Status                  | Meaning                          |
| ----------------------- | -------------------------------- |
| `401 Unauthorized`      | Incorrect password.              |
| `403 Forbidden`         | Account is suspended.            |
| `429 Too Many Requests` | Rate limit exceeded for this IP. |

---

\subsection{Student login}

`POST /auth/login/student`

Identical to general login but rejects credentials that do not belong to a `student` role. Use this endpoint in student-facing clients to prevent admins from signing in through the student app.

**Auth required:** No

\subsubsection{Request body}

<ParamField body="username" type="string" required>
  Student roll number (e.g. `O210008`).
</ParamField>

<ParamField body="password" type="string" required>
  The student's password.
</ParamField>

\subsubsection{Example}

\begin{lstlisting}[language=bash]
curl --request POST \
  --url https://api.uniz.rguktong.in/api/v1/auth/login/student \
  --header 'Content-Type: application/json' \
  --data '{
    "username": "O210008",
    "password": "password123"
  }'

\end{lstlisting}

**200 OK**

\begin{lstlisting}[language=json]
{
  "success": true,
  "token": "eyJhbGciOiJIUzI1NiIsInR5cCI6IkpXVCJ9...",
  "role": "student",
  "username": "O210008"
}

\end{lstlisting}

\subsubsection{Error responses}

| Status                  | Meaning                                                             |
| ----------------------- | ------------------------------------------------------------------- |
| `401 Unauthorized`      | Incorrect password.                                                 |
| `403 Forbidden`         | Account is suspended, or the account belongs to a non-student role. |
| `429 Too Many Requests` | Rate limit exceeded for this IP.                                    |

---

\subsection{Admin login}

`POST /auth/login/admin`

Authenticates administrative accounts only. Rejects student credentials.

**Auth required:** No

\subsubsection{Request body}

<ParamField body="username" type="string" required>
  Admin username (e.g. `director`, `dean_cse`, `warden_male`,
  `caretaker_female`, `security_admin`).
</ParamField>

<ParamField body="password" type="string" required>
  The admin's password.
</ParamField>

\subsubsection{Example}

\begin{lstlisting}[language=bash]
curl --request POST \
  --url https://api.uniz.rguktong.in/api/v1/auth/login/admin \
  --header 'Content-Type: application/json' \
  --data '{
    "username": "director",
    "password": "director@uniz"
  }'

\end{lstlisting}

**200 OK**

\begin{lstlisting}[language=json]
{
  "success": true,
  "token": "eyJhbGciOiJIUzI1NiIsInR5cCI6IkpXVCJ9...",
  "role": "director",
  "username": "director"
}

\end{lstlisting}

\subsubsection{Error responses}

| Status                  | Meaning                                                         |
| ----------------------- | --------------------------------------------------------------- |
| `401 Unauthorized`      | Incorrect password.                                             |
| `403 Forbidden`         | Account is suspended, or the account belongs to a student role. |
| `429 Too Many Requests` | Rate limit exceeded for this IP.                                |

---

\subsection{Logout}

`POST /auth/logout`

Clears the server-side session. The client should also discard the stored JWT.

**Auth required:** No

\subsubsection{Example}

\begin{lstlisting}[language=bash]
curl --request POST \
  --url https://api.uniz.rguktong.in/api/v1/auth/logout \
  --header 'Authorization: Bearer <token>'

\end{lstlisting}

**200 OK**

\begin{lstlisting}[language=json]
{
  "success": true
}

\end{lstlisting}


\section{Deep Technical Analysis}
The underlying system architecture for this module is built upon the \textbf{UniZ Microservices Framework}. 
This design ensures that each functional unit (Auth, Profile, Academics) remains isolated, preventing cascading failures across the university ecosystem.

\subsection{Operational Hardening}
Deployments are managed via K3s (Kubernetes) and containerized using high-performance Alpine Linux Docker images. 
Resource management is critical; for instance, the documentation service itself is provisioned with 2GB of RAM to ensure the responsive rendering of complex documentation graphs and state-management components.

\subsection{Enterprise Security Topology}
Every request entering the UniZ gateway is subjected to an aggressive JWT validation layer. 
Permissions are granular, mapping directly onto the RBAC (Role-Based Access Control) matrix defined in the core system specifications.

\subsection{Data Governance}
Prisma-driven data access ensures type safety and predictable query performance. 
We utilize PostgreSQL connection pooling to handle concurrent requests from the student portal and administrative dashboards simultaneously.
